\section{Spin correlations and magnetic moments}\label{mo:section}

The magnetic estimates in this section concern the critical nearest-neighbor
Ising model with no external fields except those induced by prescribed spins.
Pins outside the stated free neighborhoods may have arbitrary signs.

\begin{lemma}[Static correlation inputs]\label{fd:static-correlations}
The following statements hold with constants depending only on the indicated
fixed bulk separations and aspect ratios.
\begin{enumerate}
\item In a square of radius \(r\), under any FK boundary wiring,
\[
 c r^{-1/8}\le
 \phi(x\longleftrightarrow\partial B_r)
 \le C r^{-1/8}
\]
for \(x\) in a fixed central subsquare.  In the plus Ising law this is
the corresponding magnetization estimate.
\item In a plus rectangle of scale \(r\) and fixed aspect ratio, for two
distinct points in separated fixed compact bulk sets,
\[
 c r^{-1/4}\le\operatorname{Cov}_+(\sigma_x,\sigma_y)
 \le C r^{-1/4}.
\]
\item The infinite-volume zero-field law on \(\mathbb Z^2\) is unique and
spin-flip symmetric.  For a lattice-dependent constant \(c_\sigma>0\),
\[
 \mathbb E(\sigma_x\sigma_y)
   =c_\sigma |x-y|^{-1/4}(1+o(1))
       \quad (|x-y|\longrightarrow\infty),
 \qquad
 0\le\mathbb E(\sigma_x\sigma_y)
       \le C(1+|x-y|)^{-1/4}.
\]
The asymptotic is uniform in direction.
\end{enumerate}
\end{lemma}

\begin{proof}
The centered square assertion is
\cite[arXiv version~1, Lemma~26]{DHN2011}.
For another bulk point, take a square of radius \(c'r\) about that
point contained in the original square.  Its free-boundary arm lower
bound remains a lower bound for the arm inside the larger domain.
Require in addition an open circuit in its outer half-annulus and
a rectangle crossing from its inner half to the original boundary.
The arm, circuit, and crossing then join.  RSW gives the latter two
events jointly a fixed positive probability, and FK association
multiplies this by the arm lower bound.  The upper bound uses the
wired arm in a contained square.  This proves the stated uniformity.
The magnetization tends to zero as the boundary recedes.  The monotone
coupling of plus and minus boundary laws has disagreement probability
equal to half their mean difference at a vertex.  A union bound on any
fixed finite set therefore makes the extremal infinite-volume laws equal,
proving uniqueness and symmetry.

For the second assertion, apply bulk one- and two-spin convergence
\cite[Theorem~1.3]{CHI2021} in the fixed rescaled rectangle.  We spell out
strict positivity of the limiting connected function.  In the upper
half-plane let
\(q=|z-z'|/|z-\overline{z'}|\in(0,1)\).
The plus one- and two-spin formulae
\cite[Theorem~7.1, Equation~(7.3)]{CHI2021} give
\[
 \frac{\langle\sigma_z\sigma_{z'}\rangle_+}
      {\langle\sigma_z\rangle_+\langle\sigma_{z'}\rangle_+}
 =\left(\frac{q^{1/2}+q^{-1/2}}2\right)^{1/2}>1.
\]
Conformal covariance transfers this strict inequality to a rectangle.
Continuity and compactness give a positive minimum for the connected
function on separated compact bulk sets.  Uniform bulk convergence gives
the discrete lower bound for sufficiently large \(r\); finite energy
absorbs the finitely many smaller scales.  The upper bound also follows
from convergence, or from two disjoint wired arm tests.

For completeness, put \(R=|x-y|\) and \(u=(y-x)/R\), and bracket the
plane two-point function between free and plus laws in a square centered
at \(x\) of radius \(TR\).  Fix \(T>2\) first and take \(R\to\infty\).
Theorem~1.3 of \cite{CHI2021} gives the continuum limits
\(F_T^b(0,u)\), for \(b\in\{\mathrm{free},+\}\), uniformly for
\(|u|=1\).  Conformal scaling gives
\(F_T^b(0,u)=T^{-1/4}F_1^b(0,u/T)\).
The normalization and two-spin fusion in
\cite[Proposition~5.3, Theorem~6.3, Equation~(6.12)]{CHI2021} imply
\(F_1^b(0,u/T)=T^{1/4}(1+o_T(1))\) for both boundary conditions.
Remark~6.4 there makes the error uniform in direction: the colliding
pair remains in the bulk and there are no other insertions.  Letting
\(T\to\infty\) after the fixed-\(T\) mesh limit proves the plane
asymptotic with the common lattice conversion constant \(c_\sigma\),
uniformly in direction.  Finally, connection of \(x\) to \(y\) forces arms in the two
disjoint boxes of radius \(\lfloor|x-y|/4\rfloor\).  Conditioning on all
other edges and using the wired upper arm bound proves the uniform upper
bound.  Nonnegativity follows from the FK representation.
\end{proof}

\begin{lemma}[Product arm bound]\label{fd:fk-product}
Let \(Q\) be a lattice square with an ordinary FK bond law under
arbitrary boundary wiring.  After sampling the bonds, assign each
cluster touching the boundary any sign, and give every other cluster
an independent fair sign.  Conditional on the bonds, all nonboundary
cluster signs are mutually independent and independent of the entire
collection of boundary-cluster signs.  Boundary-cluster signs may have
any joint law conditional on the bonds; the bonds are not further conditioned on compatibility
with prescribed mixed signs.  Suppose
\(x_1,\ldots,x_n\) lie a distance at least \(cs\) from \(\partial Q\)
and in a set of diameter at most \(As\), where \(n\ge2\).
For
\[
 r_i=\max\left(1,\left\lfloor
       \min\left(\tfrac{cs}{4},
             \tfrac14\min_{j\ne i}\|x_i-x_j\|_\infty\right)
                         \right\rfloor\right),
\]
one has
\[
 \left|\mathbb E\prod_{i=1}^n\sigma_{x_i}\right|
 \le C_{n,c}\prod_{i=1}^n r_i^{-1/8}.
\]
The same bound holds conditionally on arbitrary bonds exterior to \(Q\),
with the induced wiring, before cluster colors are sampled.
\end{lemma}

\begin{proof}
First cancel insertions that occur an even number of times at one vertex.
At a vertex occurring repeatedly, the original \(r_i\)'s are one, so
discarding those factors costs nothing.  All remaining radii are the
original radii; none is recomputed after cancellation.  Each remaining singleton whose
ball has radius larger than one must have an open arm to that ball's
boundary if the conditional cluster-sign average is nonzero: otherwise
its free cluster contains no other odd insertion and has mean sign zero.
A cluster reaching a prescribed-sign boundary also requires this arm.
The balls of radius greater than one are disjoint, and are contained in
\(Q\).  Their arm probabilities, conditionally on all exterior edges,
are at most their fully wired probabilities.  Lemma~\ref{fd:static-correlations}
bounds each by \(Cr_i^{-1/8}\).  Sequential conditioning multiplies
the bounds.  Absolute cluster-sign averages are at most one.  Radius-one
balls and parity cancellations can only remove requirements, proving the
claim.  Exterior bond conditioning supplies only boundary wiring,
so the same proof applies conditionally.  Prescribed mixed-spin
boundaries are handled by the comparison in Lemma~\ref{fd:moments}.
\end{proof}

\begin{lemma}[Buffered line and area moments]\label{fd:moments}
Fix \(A,c>0\), an integer \(p\ge1\), and \(\nu\in\{1,2\}\).
Let \(S\) be a set of lattice points of diameter at most \(As\) such that
\[
 |S\cap B_t(x)|\le A(t+1)^\nu\quad(t\ge0),
 \qquad |S|\le A s^\nu.
\]
Assume the Ising law has a free square-lattice neighborhood of width
\(cs\) about \(S\), with arbitrary pins beyond that neighborhood.
Then, for weights \(|a_x|\le1\),
\begin{equation}
 \mathbb E\left|\sum_{x\in S}a_x\sigma_x\right|^{2p}
 \le C_{p,A,c}\,s^{2p(\nu-1/8)}.
 \label{fd:moment-bound}
\end{equation}
This includes a fixed number of line segments or square contours
\((\nu=1)\), and area sums \((\nu=2)\).  It applies as well in a
buffered chart of a torus or cylinder.  No pin on \(S\) is allowed in
this assertion.
\end{lemma}

\begin{proof}
Use Corollary~\ref{fd:buffered-ratio}, and a fixed finite covering if
necessary, to compare to a plus square with a proportional buffer.
It is enough to prove the bound there.  Expand the even power and take
absolute values of the spin-product expectations.  Lemma~\ref{fd:fk-product}
and the diameter bound reduce the desired sum to a constant times
\[
 \sum_{x_1,\ldots,x_{2p}\in S}
 \prod_{i=1}^{2p}
       \sum_{j\ne i}(1+\|x_i-x_j\|_\infty)^{-1/8}.
\]
Indeed the capped nearest distance in that lemma is bounded below by
a fixed multiple of the uncapped nearest distance plus one, since all
points lie within diameter \(As\).

Expand the product by choosing one neighbor \(j(i)\ne i\) for each
index.  The directed graph on the indices has outdegree one, and each
component consists of an oriented cycle of length at least two with
in-trees attached.  Put \(q=1/8\).  The counting hypothesis gives, uniformly
in \(y\),
\[
 \sum_{x\in S}(1+\|x-y\|_\infty)^{-q}
       \le C s^{\nu-q},
 \qquad
 \sum_{x\in S}(1+\|x-y\|_\infty)^{-2q}
       \le C s^{\nu-2q}.
\]
These follow by dyadic distance shells; both \(q\) and \(2q\) are less
than \(\nu\).  Summing an in-tree leaf removes one edge and contributes
\(Cs^{\nu-q}\).  Once only a cycle remains, sum one vertex incident to
two cycle edges by Cauchy--Schwarz, at cost \(Cs^{\nu-2q}\).
For a two-cycle the two edges have the same endpoints and this is the
second displayed bound.  The remaining graph is a path; remove its leaves
at cost \(Cs^{\nu-q}\) each, and its last vertex at cost \(Cs^\nu\).
A component with \(m\) vertices therefore costs at most
\(C_m s^{m(\nu-q)}\).  There are at most \((2p-1)^{2p}\) neighbor
choices, proving the moment estimate.

For a large or bent support, partition \(S\) by overlay squares of
side \(cs/16\).  At most \(C(1+A/c)^2\) such squares meet
\(S\).  In each nonempty piece choose \(z\in S\) in that square.
The piece lies in \(z+B_{cs/8}\), and \(z+B_{cs/2}\) lies in the
assumed free neighborhood because \(z\in S\).  Compare only this
local marginal to the plus law in the latter square.  The counting
hypotheses persist with constants depending only on \(A,c\).  Apply
the proved bound to each piece and use the triangle inequality in
\(L^{2p}\).  Lattice rounding at bounded \(cs\) is absorbed in the
same constants.  The same construction works in periodic charts.
\end{proof}
